% =====================================================================
%  Guion de la présentation
%  « Les transformers (type BERT) encodent-ils la syntaxe d'une langue ? »
%  Marcos Lahoz Yague --- Orange Research, 13 mai 2026
% =====================================================================
\documentclass[11pt]{article}

\usepackage[T1]{fontenc}
\usepackage[utf8]{inputenc}
\usepackage{lmodern}
\usepackage[a4paper, margin=2.3cm]{geometry}
\usepackage{xcolor}
\usepackage{titlesec}
\usepackage{enumitem}
\usepackage{microtype}
\usepackage[hidelinks]{hyperref}
\usepackage{xspace}
\usepackage{amsmath}

\definecolor{accent}{RGB}{255, 121, 0}
\definecolor{darknavy}{RGB}{26, 43, 76}
\definecolor{midgray}{RGB}{120, 130, 140}

\titleformat{\section}
  {\large\bfseries\color{darknavy}}
  {Slide~\thesection.}{0.5em}{}
\titlespacing*{\section}{0pt}{1.2em}{0.4em}

\setlength{\parindent}{0pt}
\setlength{\parskip}{0.4em}

\newcommand{\visuel}[1]{%
  {\footnotesize\color{midgray}\textit{Visuel : #1}\par\vspace{0.3em}}}
\newcommand{\transition}[1]{%
  \par\vspace{0.4em}{\footnotesize\color{accent}\textit{Transition :} #1}}
\newcommand{\duree}[1]{\hfill{\footnotesize\color{midgray}(${\sim}#1$\,s)}}
\newcommand{\mbert}{m\textsc{BERT}\xspace}

\title{\Large Guion de la présentation\\[0.3em]
       \large\textcolor{midgray}{« Les \emph{transformers} (type BERT)
       encodent-ils la syntaxe d'une langue ?»}}
\author{Marcos Lahoz Yague \textemdash{} Orange Research, 13 mai 2026}
\date{}

\begin{document}
\maketitle

\noindent\textcolor{midgray}{\footnotesize Budget total : 8 minutes
(${\sim}480$\,s, soit ${\sim}25$--$30$\,s par slide en moyenne).
Les durées indicatives sont à droite de chaque section.}

\hrule
\vspace{0.6em}

% =====================================================================
\section{Titre} \duree{20}
\visuel{page de titre, nom et affiliation.}

Bonjour, je suis Marcos Lahoz Yague. Aujourd'hui je vais vous présenter
un travail réalisé pendant mon Master 2 à Paris-Saclay, sur la
question : les modèles \emph{transformers} de type BERT,
encodent-ils la syntaxe du langage ? En huit minutes, je vais vous
montrer la méthode, les résultats principaux, et où ce travail
pourrait s'étendre vers le doctorat.

% =====================================================================
\section{Le langage est arborescent} \duree{30}
\visuel{phrase « The cat ate the fish» avec son arbre de
        dépendances.}

Pour répondre à cette question, il faut d'abord se mettre d'accord sur
ce que nous appelons « syntaxe». Dans ce travail, nous prenons
une définition précise : l'arbre de dépendances. Chaque mot d'une
phrase a exactement \emph{une} tête grammaticale --- le mot dont il
dépend. La seule exception est la racine, typiquement le verbe
principal. Les arêtes peuvent porter des étiquettes --- sujet, objet,
déterminant --- mais comme Hewitt et Manning, nous ne regardons que
la \emph{topologie}. C'est cet arbre que nous chercherons à retrouver
dans BERT.

% =====================================================================
\section{Deux métriques sur l'arbre} \duree{35}
\visuel{profondeur (badges orange) et distance (chemin orange) sur le
        même arbre.}

Cette définition nous donne deux quantités numériques naturelles. À
gauche : la \textbf{profondeur} d'un mot --- le nombre d'arêtes
jusqu'à la racine. Le verbe principal est à zéro, ses dépendants
directs à un. À droite : la \textbf{distance} entre deux mots --- le
nombre d'arêtes du chemin unique qui les relie. Par exemple,
« cat» et « fish» sont à distance deux, parce qu'on passe
par « ate». C'est sur la distance que nous concentrons nos
expériences --- elle donne $O(n^2)$ paires de supervision par phrase
contre $O(n)$ pour la profondeur, et c'est aussi sur les distances
que la sonde isométrique aura un sens géométrique direct.

% =====================================================================
\section{BERT comprend-il la syntaxe ?} \duree{30}
\visuel{liste des tâches où BERT excelle + diagramme MLM.}

BERT est un modèle \emph{transformer} préentraîné qui a marqué une
étape majeure en NLP. Il obtient d'excellents résultats sur de
nombreuses tâches : réponse aux questions, inférence textuelle,
classification, recherche web. Ce succès suggère que BERT
« comprend» le langage. Et pourtant --- c'est ce que montre
le diagramme à droite --- pendant son préentraînement, BERT ne fait
qu'\emph{une} chose : prédire des tokens masqués. Aucun arbre, aucune
étiquette grammaticale, aucune tête, aucun rôle. La question est
donc : si la syntaxe est là, où est-elle ?

% =====================================================================
\section{Qu'est-ce qu'une sonde ?} \duree{40}
\visuel{schéma de la sonde au-dessus du modèle figé.}

La méthode pour répondre s'appelle le \emph{probing}. Le principe :
on gèle les paramètres du modèle, on extrait les vecteurs $h_i$
d'une couche cachée, et au-dessus on entraîne un \emph{petit} modèle
supervisé --- typiquement une simple projection linéaire --- à
prédire une propriété linguistique annotée. Pourquoi un modèle aussi
petit ? Parce que si une projection linéaire suffit, c'est que la
propriété est \emph{déjà} encodée géométriquement dans les vecteurs.
Une sonde trop puissante apprendrait elle-même la structure ---
même à partir de vecteurs aléatoires --- et on ne saurait plus si
l'information vient du modèle ou de la sonde.

% =====================================================================
\section{La sonde structurelle de Hewitt \& Manning (2019)} \duree{45}
\visuel{formule $\|B(h_i - h_j)\|^2 \approx d_T(w_i, w_j)$ +
        flèche $B$ des vecteurs aux points projetés.}

En 2019, Hewitt et Manning proposent la \emph{sonde structurelle}.
L'idée : apprendre une projection linéaire $B$ telle que la distance
euclidienne au carré entre deux vecteurs projetés approxime la
distance dans l'arbre de dépendances. Visuellement, à droite : on
prend les vecteurs bruts de BERT (les points bleus en haut), on les
projette par $B$, et on obtient les points oranges en bas, disposés
de manière à ce que leurs distances euclidiennes reflètent l'arbre.
La sonde n'a accès qu'aux vecteurs --- jamais à la phrase, jamais à
l'arbre directement. Si cette projection linéaire fonctionne, alors
la syntaxe vit géométriquement dans l'espace des représentations.

% =====================================================================
\section{Comment évaluer la sonde ? Deux métriques} \duree{45}
\visuel{scatter plot Spearman + arbre UUAS avec une arête erronée.}

Pour mesurer la qualité de la sonde, deux métriques. À gauche, la
\textbf{corrélation de Spearman} : on regarde si les distances
prédites sont dans le \emph{bon ordre} que les distances vraies.
C'est une corrélation de rang, calculée sur toutes les paires de
mots d'une phrase. On rapporte $\rho_{5\text{-}50}$, moyenne sur les
phrases de longueur 5 à 50 mots. À droite, l'\textbf{UUAS} : à
partir des distances prédites, on construit un arbre couvrant
minimum, et on compte la fraction d'arêtes correctes --- non
étiquetées, non orientées. Dans l'exemple, trois sur quatre arêtes
sont correctes, soit UUAS $= 0{,}75$.

% =====================================================================
\section{Résultat sur l'anglais} \duree{30}
\visuel{graphique de barres : Baseline vs BERT-base, UUAS et Spearman.}

Voici le résultat principal de Hewitt et Manning sur l'anglais. À
droite, BERT-base à la couche 7 : UUAS $= 0{,}80$, Spearman
$= 0{,}85$. À gauche, le baseline non contextualisé : $0{,}49$ et
$0{,}50$. Plus 31 points en UUAS, plus 35 en Spearman --- sans aucun
entraînement syntaxique. Performance proche des analyseurs
syntaxiques entièrement supervisés sur PTB. La syntaxe vit donc bien
dans la géométrie de BERT.

% =====================================================================
\section{Ce que nous ajoutons} \duree{25}
\visuel{deux questions ouvertes en colonnes.}

Hewitt et Manning ont répondu à la question pour BERT en anglais.
Nous prolongeons leur travail avec deux questions ouvertes.
Premièrement : la même géométrie syntaxique existe-t-elle dans
\mbert{}, pour des langues autres que l'anglais ? Concrètement, ça
marche-t-il sur l'espagnol ? Deuxièmement : la syntaxe est-elle une
forme \emph{rigide} dans l'espace, ou faut-il l'\emph{amplifier}
par des étirements non uniformes ? Pour ces deux questions, deux
contributions : une réplication couche par couche sur AnCora, et
une \emph{sonde isométrique}.

% =====================================================================
\section{Données : UD\_Spanish-AnCora} \duree{25}
\visuel{partitions du treebank, encodeur, hyperparamètres de la sonde.}

Pour la première question, nous utilisons le treebank
\textsc{ud\_spanish-ancora} --- la partie espagnole d'AnCora dans
Universal Dependencies. Quatorze mille phrases d'entraînement, mille
six cents en dev, mille sept cents en test. L'encodeur est \mbert{},
\texttt{bert-base-multilingual-cased}, couvrant 104 langues,
complètement figé. La sonde a un rang 128, perte $\ell_1$, 30 époques
avec early stopping. L'espagnol amène deux complications absentes du
PTB anglais : les \emph{contractions} comme « del = de + el»,
et les sujets \emph{pro-drop} sous forme de n\oe{}uds vides --- il
faut bien les filtrer pour ne pas polluer la sonde.

% =====================================================================
\section{Résultat 1 : layer sweep sur l'espagnol} \duree{40}
\visuel{courbe en cloche sur les 13 couches de mBERT.}

Nous balayons les 13 couches de \mbert{} et entraînons une sonde par
couche. Le résultat : une courbe en cloche qui reproduit exactement
la forme observée par Hewitt et Manning sur l'anglais. Spearman
pique à $0{,}86$ à la couche 7, UUAS à $0{,}76$ à la même couche.
Donc la syntaxe espagnole se concentre au milieu du réseau, comme
en anglais.

% =====================================================================
\section{Transfert cross-lingue} \duree{40}
\visuel{tableau comparatif H\&M vs nous.}

Comparons directement à Hewitt et Manning. Eux, sur PTB anglais :
Spearman ${\approx} 0{,}85$, UUAS ${\approx} 0{,}80$. Nous, sur
AnCora espagnol via \mbert{} : Spearman $= 0{,}86$, UUAS $= 0{,}76$.
La même couche pique. Spearman est même légèrement plus haut. UUAS
n'est que $4$ points en dessous, malgré la dilution de \mbert{} sur
104 langues. C'est une forte évidence pour un \emph{sous-espace
syntaxique partagé} entre les langues, en accord avec Chi et al.\
(ACL 2020).

% =====================================================================
\section{La sonde isométrique} \duree{45}
\visuel{formule $B B^\top = I$ + interprétation rigide vs étirable.}

Passons maintenant à la deuxième question. La sonde linéaire fait
deux choses entremêlées : elle \emph{tourne} et elle \emph{étire}
l'espace. Notre extension isole ces deux opérations en contraignant
la projection à être orthogonale : $B B^\top = I$, c'est-à-dire que
$B$ est une rotation pure. Si la syntaxe est une forme rigide dans
l'espace \mbert{}, l'isométrique doit égaler la sonde linéaire. Si
la syntaxe nécessite un étirement anisotrope, l'isométrique doit
s'effondrer. C'est un test propre, par construction.

% =====================================================================
\section{Deux scénarios} \duree{25}
\visuel{illustration des deux hypothèses --- rigide vs étirable.}

Les deux scénarios qu'on cherche à distinguer. À gauche : si l'arbre
est rigide dans le sous-espace, l'isométrique conserve les distances
et reproduit le résultat de la sonde linéaire. À droite : si l'arbre
est dans une direction \emph{étirée} de faible variance, la rotation
pure ne peut pas le récupérer.

% =====================================================================
\section{Résultat 2 : l'effet Pancake} \duree{50}
\visuel{tableau Linéaire vs Isométrique + interprétation Pancake.}

Le résultat est sans appel. À la couche 7, la sonde linéaire :
UUAS $= 0{,}76$, Spearman $= 0{,}86$, nombre de condition
$\kappa(B) = 3{,}74$ --- la sonde étire l'espace par un facteur
${\sim}4$ entre ses directions principales. La sonde isométrique,
contrainte à être une rotation pure, s'effondre :
UUAS $= 0{,}27$, Spearman $= 0{,}35$ --- soit \textbf{$-64\,\%$ en UUAS}. Interprétation : l'arbre syntaxique espagnol est aplati dans
des directions de \emph{faible variance} de l'espace \mbert{}. La
sonde doit amplifier ces directions par environ $4$ pour faire
ressortir la structure. C'est ce que nous appelons l'effet
\emph{Pancake} : la syntaxe vit dans une crêpe, pas dans une boule.

% =====================================================================
\section{Limitations} \duree{25}
\visuel{liste des limitations.}

Quelques limites honnêtes. Tous les chiffres sont sur \emph{une
seule} graine ; il faudrait répéter avec au moins 5 seeds et reporter
moyenne $\pm$ écart-type. Le test set n'a pas été touché, par
discipline. Une tâche de contrôle à la Hewitt-Liang aborderait la
question de la capacité intrinsèque de la sonde. Et il faudrait
comparer contre des encodeurs espagnols monolingues comme BETO,
pour quantifier le coût de la dilution multilingue de \mbert{}.

% =====================================================================
\section{Du probing syntaxique au probing multimodal} \duree{40}
\visuel{tableau de correspondances syntaxe → multimodal.}

Comment ce travail se relie-t-il au doctorat sur le raisonnement
multimodal sur vidéo ? La méthodologie --- sonde plus diagnostic
géométrique --- se généralise directement. Mêmes questions, autres
propriétés : où vit chaque modalité dans un MLLM ? Sont-elles
rigidement préservées, c'est-à-dire $\kappa$ proche de 1 quand on
isole chacune ? À quelle couche se cristallise leur interaction,
et observe-t-on un effondrement modal analogue à l'effet Pancake ?
Cela connecte avec les travaux récents sur le \emph{modality
collapse} (Park et al., AAAI 2025 ; Wang et al., CVPR 2025). Même
méthodologie, dimension supérieure.

% =====================================================================
\section{À retenir} \duree{30}
\visuel{trois points résumés.}

Trois choses à retenir.

\textbf{Premièrement}, nous avons répliqué la sonde structurelle
en espagnol --- UUAS $0{,}76$ à la couche 7, à seulement $4$ points
du pic anglais. Forte évidence d'un sous-espace syntaxique
cross-lingue.

\textbf{Deuxièmement}, nous avons documenté trois pièges
méthodologiques de réplication qui empoisonnaient silencieusement le
pipeline (n\oe{}uds vides, ponctuation espagnole, indexation des
couches) : leur correction donne $+51\,\%$ d'UUAS.

\textbf{Troisièmement}, nous avons introduit la \emph{sonde
isométrique}. La chute de $-64\,\%$ d'UUAS montre que la géométrie
syntaxique espagnole est \emph{étirable, pas rigide} : effet
« Pancake», $\kappa(B) = 3{,}74$.

\medskip

Merci. Je reste à votre disposition pour les questions.

\end{document}
